% !TeX root = ../../main.tex
% Jablonski diagram of organic scintillation: singlet ladder on the left,
% triplet on the right, one arrow per process; legend and a ternary-plastic
% inset below. \input by chapters/02-literature-review.tex. Companion to
% scintillation-process.tex and drawn to its palette and weights.
%
% Labels sit beside their own arrow, never across the gap between the singlet
% and triplet columns: that gap (x = 6.6 to 8.6) belongs to the intersystem
% crossing. The delayed-fluorescence return goes over the top at y = 5.0 so it
% cannot run along the S1 line that intersystem crossing leaves from. Every
% prose block has a text width, so its box is fixed rather than set by the
% wording. \orgfit guards the total width, which the right-hand labels set.
%
% Needs arrows.meta and decorations.pathmorphing, both loaded in main.tex.

\providecolor{orgflu}{HTML}{378ADD}
\providecolor{orgphos}{HTML}{D4537E}
\providecolor{orgdelay}{HTML}{EF9F27}
\providecolor{orgexcite}{HTML}{E24B4A}
\providecolor{orghot}{HTML}{BA7517}
\providecolor{orginert}{HTML}{888780}
\providecolor{orgboxfill}{HTML}{FAEEDA}

\providecommand{\orgfit}{%
	\path let \p1=(current bounding box.east), \p2=(current bounding box.west) in
		\pgfextra{%
			\ifdim\dimexpr\x1-\x2\relax>\textwidth
				\PackageError{organic-scintillation}{Figure is wider than the text block}%
					{It measures \the\dimexpr\x1-\x2\relax\space against a measure of
					\the\textwidth. Shorten a label or move it inside the diagram.}%
			\fi};}

\begin{tikzpicture}[
	x=1cm, y=1cm,
	ann/.style   = {font=\footnotesize, align=center},
	side/.style  = {ann, align=left},
	title/.style = {font=\small},
	lvl/.style   = {line width=0.8pt},
	vib/.style   = {orginert!50, line width=0.3pt},
	flow/.style  = {-{Stealth[length=4pt]}, line width=0.5pt},
	relax/.style = {-{Stealth[length=3pt]}, orginert, line width=0.4pt, decorate,
		decoration={snake, amplitude=0.35mm, segment length=1.4mm, post length=1.2mm}},
	wavy/.style  = {-{Stealth[length=4pt]}, line width=0.5pt, decorate,
		decoration={snake, amplitude=0.5mm, segment length=2mm, post length=1.6mm}},
	box/.style   = {draw=orgdelay, fill=orgboxfill, line width=0.4pt, rounded corners=1pt,
		minimum height=0.8cm, inner xsep=3pt, font=\footnotesize, align=center,
		text=orgdelay!55!black, text width=2.8cm},
	]

% ---------- levels: S0, S1, S2 at left, T1 at right ----------------
% {y}/{label}: a thick line with three vibrational sublevels above it
\foreach \y/\l in {0.4/$S_0$, 4.2/$S_1$, 6.6/$S_2$} {
	\draw[lvl] (1.2,\y) -- (6.6,\y);
	\foreach \d in {0.22,0.44,0.66} \draw[vib] (1.2,\y+\d) -- (6.6,\y+\d);
	\node[ann, anchor=east] at (0.95,\y) {\l};
}
\draw[lvl] (8.6,3.0) -- (12.0,3.0);
\foreach \d in {0.22,0.44,0.66} \draw[vib] (8.6,3.0+\d) -- (12.0,3.0+\d);
\node[ann, anchor=east] at (8.35,3.0) {$T_1$};
\draw[lvl] (8.6,0.4) -- (12.0,0.4);

\draw[flow, line width=0.3pt] (0.4,0.4) -- (0.4,7.3);
\node[ann] at (0.4,7.55) {$E$};

\node[title] at (3.9,8.2)  {Singlet states};
\node[ann]   at (3.9,7.75) {spins paired: emission allowed};
\node[title] at (10.3,8.2) {Triplet state};
\node[ann]   at (10.3,7.75) {spins parallel: emission forbidden};

% ---------- excitation and the two singlet relaxations -------------
\draw[flow, orgexcite, line width=0.6pt] (1.8,0.4) -- (1.8,7.26);
\node[side, anchor=west, text width=3.1cm] at (2.0,1.7)
	{Ionising particle\\[-1pt] excites $\pi$ electrons,\\[-1pt] $\sim$\qty{e-15}{\second}};

\draw[relax] (2.6,7.26) -- (2.6,6.64);
\draw[flow, orginert, dashed] (3.6,6.6) -- (3.6,4.88);
\draw[relax] (3.6,4.86) -- (3.6,4.24);
\node[side, anchor=west, text width=3.3cm] at (3.8,5.75)
	{Internal conversion\\[-1pt] non-radiative, $\sim$\qty{e-12}{\second}};

\draw[flow, orgflu] (5.4,4.2) -- (5.4,0.86);
\draw[relax] (5.4,0.84) -- (5.4,0.44);
\node[side, anchor=east, text width=3.1cm, align=right] at (5.2,2.95)
	{Fluorescence\\[-1pt] \qtyrange[range-phrase=--, range-units=single]{1}{3}{\nano\second}, prompt\\[-1pt] $\sim$\qty{420}{\nano\metre} in plastics};
\draw[wavy, orghot] (5.55,1.55) -- (6.5,1.2);

% ---------- the triplet detour --------------------------------------
\draw[flow, orginert, dashed] (6.6,4.2) -- (8.6,3.68);
\draw[relax] (8.9,3.66) -- (8.9,3.04);
\node[ann, text width=2.2cm] at (7.4,2.05) {Intersystem crossing};

\draw[flow, orgphos, dotted, line width=0.6pt] (10.3,3.0) -- (10.3,0.86);
\draw[relax] (10.3,0.84) -- (10.3,0.44);
\node[side, anchor=west, text width=3.2cm] at (10.5,1.6)
	{Phosphorescence\\[-1pt] spin-forbidden, $\sim$ms,\\[-1pt] little light};

% over the top: the return to S1 must not share the S1 line with the
% intersystem-crossing arrow
\draw[flow, orgdelay] (12.0,3.0) -- (12.6,3.0) -- (12.6,5.0) -- (6.0,5.0) -- (6.0,4.3);
\node[side, anchor=west, text width=2.85cm] at (12.7,4.0)
	{Delayed\\[-1pt] fluorescence:\\[-1pt] $T_1 + T_1 \rightarrow S_1 + S_0$\\[-1pt] refills $S_1$, so the\\[-1pt] same photon comes\\[-1pt] hundreds of ns late};

% ---------- footnote ------------------------------------------------
\node[ann, anchor=north west, text width=14.6cm, align=left] at (0.5,-0.45)
	{Both emissions terminate above the $S_0$ ground vibrational level, so the emitted photon is red-shifted from the absorption and the material is transparent to its own light.};

% ---------- legend ---------------------------------------------------
\draw[flow, orgflu] (0.6,-1.7) -- (1.3,-1.7);
\node[ann, anchor=west] at (1.4,-1.7) {Radiative, spin-allowed};
\draw[flow, orgphos, dotted, line width=0.6pt] (5.4,-1.7) -- (6.1,-1.7);
\node[ann, anchor=west] at (6.2,-1.7) {Radiative, spin-forbidden};
\draw[flow, orginert, dashed] (10.3,-1.7) -- (11.0,-1.7);
\node[ann, anchor=west] at (11.1,-1.7) {Non-radiative};
\draw[relax] (13.6,-1.7) -- (14.3,-1.7);
\node[ann, anchor=west] at (14.4,-1.7) {Heat};

% ---------- inset: the ternary plastic ------------------------------
\node[ann, anchor=west, align=left] at (0.5,-2.35) {In a plastic the levels above belong to the fluor, not the bulk:};
\node[box] (base) at (2.0,-3.4)  {Polymer base (PVT)\\[-1pt] absorbs the deposit};
\node[box] (fluor) at (7.6,-3.4) {Primary fluor,\\[-1pt] $\sim$1\,\% by mass:\\[-1pt] fluoresces as above};
\node[box] (wls) at (13.2,-3.4)  {Wavelength shifter\\[-1pt] re-emits to match the sensor};
\draw[flow, orginert, dashed] (base.east) -- (fluor.west)
	node[ann, midway, above=3pt] {F\"orster transfer};
\draw[wavy, orghot] (fluor.east) -- (wls.west)
	node[ann, midway, above=3pt] {UV photon};

\orgfit
\end{tikzpicture}
